\section*{Changes from earlier versions, including four corrected statements}
\addcontentsline{toc}{section}{Changes from earlier versions}

A reader who cited version~1 or version~2 of this preprint should read this section. Four statements
across those two versions are wrong as written; they are listed first, most recent first, and
separately from the new work.

\paragraph{Version 3: one claim asserted the converse of a sufficient condition.} Versions 1 and 2
reported that every released dictionary trips the failure threshold of
Corollary~\ref{cor:failthresh} earlier than its shape-matched i.i.d.\ control, and glossed that in
the abstract as: \emph{where a random code of the same shape is still safe, a released dictionary
already has a feature no affine rule recovers}. The measurement is correct --- all
\SaeFailEarlier{} of \SaeDicts{} comparisons, verified twice against the stored leverages --- and
the second clause follows from it. The first does not. The corollary is \emph{sufficient} for
non-separability, and Proposition~\ref{prop:oneway} of the same manuscript shows its converse
fails, so a threshold that has not been reached says nothing about recoverability. Writing $K$ for
an actual first non-separable sparsity and $S$ for a threshold, the measurements give
$K_{\text{dict}}\le S_{\text{dict}}$, $K_{\text{ref}}\le S_{\text{ref}}$ and
$S_{\text{dict}}<S_{\text{ref}}$, which leaves the order of $K_{\text{dict}}$ and $K_{\text{ref}}$
open. The claim stood in the abstract, in the contributions and in the discussion of both earlier
versions; all three are rewritten here as a comparison of certificates, and the reading is retired
as item~9 of Appendix~\ref{app:withdrawn}.

Three facts that bear on how much the comparison carries are now reported rather than omitted: the
controls are \SaeFailRefs{} draws shared across \SaeDicts{} dictionaries rather than \SaeDicts{}
independent ones; the range of \SaeFailGapMin{} to \SaeFailGapMax{} levels is given with its median
of \SaeFailGapMedian{} and with the control thresholds, so its two ends are not read against each
other; and at the operating sparsities the same criterion rules out every feature of
\SaeFailCtlAllRuledOut{} of the \SaeFailRefs{} controls. The direction survives at the $q=0.001$,
$0.01$ and $0.05$ leverage quantiles as well as at the minimum, so the sign does not rest on one
order statistic; the magnitude does, falling to a median of \SaeFailQFiveGapMedian{} at $q=0.05$.

\paragraph{Version 2: three statements that were too strong, and are corrected here.} An external adversarial
audit found, and we verified independently against the implementation, that:

\begin{itemize}[leftmargin=*,itemsep=3pt]

\item Theorem~\ref{thm:arrow} and Corollary~\ref{cor:failthresh} were stated for all $h_i<1$, while
their shared proof establishes them only for $h_i\le 1/2$. Explicit unit-norm counterexamples at
$d=2$, $F=3$ give $\kappa=\infty$ against a finite bound. Both now carry the $h_i\le 1/2$
hypothesis, $\kappa_i$ is defined as $+\infty$ when the linear programme is infeasible, and the
counterexamples are regression tests in the released code
(\code{tests/test\_arrow\_counterexamples.py}).

\item A proposition concluded that a feature with $\kappa_i=1$ remains separable at $s=1$.
Theorem~\ref{thm:frontier} requires $\kappa_i>s$, so the correct conclusion is that such a feature is
separable at no positive sparsity. Corrected.

\item Theorem~\ref{thm:codefloor} claimed that the row-wise optimum coincides with the algebraic
pseudoinverse. It is the \emph{gain-calibrated} pseudoinverse $D_h^{-1}\Phi^{+}$. The clause is
removed.

\end{itemize}

No measurement in v1 or here depends on the invalid region. The corollary is only ever applied at
$\min_i h_i$, and the largest $\min_i h_i$ over all \KdThreeModels{} models of the main grid is
\KdThreeMaxLevMin, inside $h_i\le 1/2$. The tests that reported no violations sampled only that
region, which is a gap in the tests rather than a false test, and is why the error survived to
publication.

\paragraph{Two readings withdrawn.} A defect in how a shared decision threshold was selected --- the
search never scored the default it was meant to improve on --- invalidated two empirical claims made
in v1: that the two pre-registered estimands disagree, and that an affine probe of the preactivation
beats the network by a full sparsity level. Refitted correctly, the estimands agree at every width
and the pre-ReLU advantage is at most a quarter of a level on trained codes, reversing at the largest
width. Appendix~\ref{app:withdrawn} records these and every other reading this work retired, with the
cause of each.

\paragraph{What is new.} The theory is unchanged apart from the corrections above. The empirical
content is substantially new: the leverage identity of Proposition~\ref{prop:leverage} is put to
predictive use rather than stated in passing; \SaeDicts{} released decoder matrices are measured
against shape-matched controls, with a trained-versus-randomly-initialised pair locating the cause;
the controlled campaigns are extended to four widths, two feature loads and two training sparsities;
a frozen-code arm separates what training the code contributes from what training the readout
contributes; and the ceiling on the decoder comparison is measured rather than argued. Every measured
number is generated from the released results by \code{scripts/make\_numbers.py} and checked against
this manuscript by \code{scripts/check\_manuscript\_numbers.py}.
